\documentclass[12pt,a4paper]{article}
\usepackage[a4paper,left=21mm,right=21mm,top=22mm,bottom=24mm]{geometry}
\usepackage[T2A]{fontenc}
\usepackage[utf8]{inputenc}
\usepackage[russian]{babel}
\usepackage{amsmath,amssymb,mathtools}
\usepackage{array,booktabs,longtable}
\usepackage{xcolor}
\usepackage{hyperref}
\usepackage{tikz}
\hypersetup{colorlinks=true,linkcolor=blue!55!black,urlcolor=blue!55!black}
\setlength{\parindent}{0pt}
\setlength{\parskip}{5pt}
\setlength{\emergencystretch}{2em}
\renewcommand{\arraystretch}{1.25}
\newcommand{\heading}[1]{\par\vspace{4pt}\textbf{#1}\par}
\newcommand{\err}{\textcolor{red}{\textbf{Ошибка:}}}
\newcommand{\idea}{\textcolor{blue!60!black}{\textbf{Ключевая идея:}}}
\newcommand{\result}{\textcolor{teal!60!black}{\textbf{Итог:}}}
\newcommand{\nZ}{n\in\mathbb Z}
\begin{document}
\thispagestyle{empty}
\vspace*{0.29\textheight}
\begin{center}
{\LARGE\bfseries Ревью домашней работы\par}
\vspace{25mm}
{\large Автор: Лёвин Артём Александрович\par}
\vspace{9mm}
{\large 19 октября 2026 года\par}
\end{center}
\vfill
\begin{center}
\begin{tikzpicture}[x=1cm,y=1cm]
\draw[blue!35!black,opacity=.68,line width=.65pt]
(0,0) .. controls (1.6,.45) and (3.6,-.45) .. (5.2,0)
.. controls (6.7,.38) and (8.2,-.3) .. (9.6,.05);
\end{tikzpicture}
\end{center}
\newpage
\section*{Сводная таблица выявленных ошибок}
В документ включены только задания, при решении которых найдены ошибки в рассуждениях или вычислениях. Условия приведены по исходному листу заданий.
\small
\begin{longtable}{@{}p{.82cm}p{2.35cm}p{4.55cm}p{3.0cm}p{3.55cm}@{}}
\toprule
№ & Тема & Суть ошибки & Ваш ответ & Правильный ответ\\
\midrule
\endfirsthead
\toprule
№ & Тема & Суть ошибки & Ваш ответ & Правильный ответ\\
\midrule
\endhead
\hyperref[task:2]{2} & Кредит: выплаты & Восемь процентов подставлены как $0{,}8$ вместо $0{,}08$. & Всего $4500$, переплата $3000$ & Всего $1800$, переплата $300$ (тыс. руб.)\\
\hyperref[task:4]{4} & Процентная ставка & Пятнадцать процентов переплаты приняты за число $15$ вместо $0{,}15$. & $200\%$ в месяц & $2\%$ в месяц\\
\hyperref[task:8]{8} & Понижение степени & В условии записано $3/5$ вместо $3/4$; в формуле для $\cos^2$ ошибочен знак. Итоговый ответ совпал с верным. & $x=\pm\frac\pi{18}$\newline $+\frac{\pi n}{3}$ & Тот же ответ; ход решения требует исправления\\
\hyperref[task:9]{9} & Степени, тригонометрия & В формуле для $\cos^2x$ неверен знак; аргумент косинуса затем поделён на два лишний раз. & $x=\frac\pi8+\frac{\pi n}{4}$ & $x=\frac\pi4+\frac{\pi n}{2}$\\
\hyperref[task:10]{10} & Синус и косинус & Из $1-\tg(x/2)=0$ неверно получено $\tg(x/2)=-1$. Итог верен. & $x=2\pi n$ или\newline $x=\frac\pi2+2\pi n$ & Тот же ответ; исправить переход\\
\bottomrule
\end{longtable}
\normalsize
Здесь и далее $n\in\mathbb Z$.
\newpage
\thispagestyle{empty}
\vspace*{\fill}
\begin{center}{\LARGE\bfseries Разбор ошибок}\end{center}
\vspace*{\fill}
\newpage
\phantomsection\label{task:2}
\section*{Задание № 2. Общая сумма выплат и переплата по кредиту}
\heading{Текст задания}
Кредит $1500$ тыс. руб. взят на четыре года под $8\%$ годовых по дифференцированной схеме. Найдите общую сумму выплат и переплату.
\heading{Ваше решение}
Вы правильно выбрали формулу суммы выплат при дифференцированном погашении:
\[
S_{\text{общ}}=S+\frac{Sp(N+1)}2,
\qquad p=\frac r{100}.
\]
Однако при вычислении переплаты записано:
\[
\frac{1500\cdot0{,}8\cdot5}{2}
=750\cdot0{,}8\cdot5=3000,
\]
а затем получена общая сумма $1500+3000=4500$.
\heading{Ваш ответ}
Переплата: $3000$. Общая сумма выплат: $4500$.
\heading{Ошибка}
\err{} Перепутана десятичная запись процентной ставки.
\heading{Где именно возникает ошибка}
До подстановки чисел рассуждение верное: общая сумма выплат равна основному долгу плюс все проценты. Первый неверный переход --- подстановка $0{,}8$ вместо $0{,}08$ при ставке $8\%$:
\[
8\%=\frac8{100}=0{,}08\ne0{,}8.
\]
\heading{Почему этот шаг неверен}
Процент означает сотую долю. Число $0{,}8$ составляет восемь десятых, то есть $80\%$, и в десять раз превышает указанную ставку. Поэтому и переплата, и сумма выплат в Ваших вычислениях завышены. При этом сама формула была выбрана верно; исправлять требуется числовую подстановку.
\heading{Ключевая идея}
\idea{} Сначала переводим ставку в десятичную дробь, а затем подставляем её в формулу. При погашении основного долга равными частями остаток долга уменьшается каждый период, поэтому сумма начисленных процентов образует арифметическую прогрессию.
\heading{Исправление от последнего верного шага}
Возвращаемся к выбранной Вами формуле и ставим $p=0{,}08$:
\[
\frac{Sp(N+1)}2=\frac{1500\cdot0{,}08\cdot5}{2}=300.
\]
\heading{Полное правильное решение}
Ежепериодная часть основного долга равна
\[
\frac{1500}{4}=375.
\]
Перед платежами остаток долга составляет соответственно $1500$, $1125$, $750$ и $375$. Начисленные проценты равны
\[
1500\cdot0{,}08=120,\quad1125\cdot0{,}08=90,
\quad750\cdot0{,}08=60,\quad375\cdot0{,}08=30.
\]
Сумма всех процентов:
\[
120+90+60+30=300.
\]
Тело кредита в сумме равно $1500$, поэтому
\[
S_{\text{общ}}=1500+300=1800.
\]
То же значение даёт сокращённая формула:
\[
S_{\text{общ}}
=1500+\frac{1500\cdot0{,}08\cdot(4+1)}2=1800.
\]
\heading{Проверка}
Полные платежи равны $375+120=495$, $375+90=465$, $375+60=435$, $375+30=405$. Их сумма $495+465+435+405=1800$, что совпадает с вычислением по формуле.
\heading{Итог}
\result{} Переплата равна $300$ тыс. руб., общая сумма выплат равна $1800$ тыс. руб.
\heading{Задача на закрепление}
Кредит $2400$ тыс. руб. погашается равными частями основного долга за три года под $6\%$ годовых. Найдите переплату и общую сумму выплат.
\newpage
\phantomsection\label{task:4}
\section*{Задание № 4. Нахождение процентной ставки}
\heading{Текст задания}
Кредит с равномерным погашением основного долга взят на $14$ месяцев. Общая сумма платежей превышает сумму кредита на $15\%$. Найдите месячную процентную ставку.
\heading{Ваше решение}
Вы использовали формулу переплаты, записали
\[
15=\frac{p(1+14)}2,\qquad30=15p,\qquad p=2,
\]
и получили
\[
\frac r{100}=2,\qquad r=200\%.
\]
\heading{Ваш ответ}
$200\%$.
\heading{Ошибка}
\err{} Относительная переплата $15\%$ подставлена как число $15$.
\heading{Где именно возникает ошибка}
Верно, что переплата как доля первоначального долга равна $p(N+1)/2$. Но здесь нужно записать
\[
\frac{15}{100}=\frac{15p}{2},
\]
а не $15=15p/2$.
\heading{Почему этот шаг неверен}
Есть три разные величины: рублёвая переплата, её доля от кредита и процентное выражение этой доли. Если переплата составляет $15\%$, доля переплаты от первоначальной суммы равна $0{,}15$. При подстановке числа $15$ доля искусственно увеличивается в сто раз. Из-за этого в сто раз увеличивается и рассчитанная процентная ставка.
\heading{Ключевая идея}
\idea{} В формулах используются доли, а не числа с обозначением процента. Если величина составляет пятнадцать процентов от целого, соответствующая доля равна пятнадцати сотым.
\heading{Исправление от последнего верного шага}
Ваша формула правильная. Исправляем только левую часть уравнения:
\[
0{,}15=\frac{15p}{2},\qquad 0{,}30=15p,\qquad p=0{,}02.
\]
Следовательно, $r=100p=2\%$.
\heading{Полное правильное решение}
Пусть первоначальный кредит равен $S$. При $14$ равных выплатах основного долга сумма остатков, на которые начисляются проценты, имеет вид
\[
S+\frac{13S}{14}+\frac{12S}{14}+\cdots+\frac S{14}
=\frac{S(14+1)}2.
\]
Отсюда общая переплата по процентам
\[
P=\frac{Sp(14+1)}2.
\]
По условию $P=0{,}15S$, поэтому
\[
0{,}15S=\frac{15Sp}2.
\]
Так как сумма кредита положительна, делим обе части на $S$ и решаем полученное линейное уравнение:
\[
0{,}15=7{,}5p,\qquad p=\frac{0{,}15}{7{,}5}=0{,}02.
\]
Переходим от доли к процентам:
\[
r=0{,}02\cdot100\%=2\%.
\]
\heading{Проверка}
Если ставка равна $2\%$, доля переплаты за четырнадцать периодов составляет
\[
\frac{0{,}02\cdot15}{2}=0{,}15=15\%,
\]
как и сказано в условии.
\heading{Итог}
\result{} Месячная процентная ставка равна $2\%$.
\heading{Задача на закрепление}
Кредит погашается равными частями основного долга за девять месяцев. Общая переплата составляет $20\%$ от первоначальной суммы кредита. Найдите месячную процентную ставку.
\newpage
\phantomsection\label{task:8}
\section*{Задание № 8. Уравнение с квадратом косинуса}
\heading{Текст задания}
Решите уравнение $\cos^2(3x)=\dfrac34$, используя понижение степени.
\heading{Ваше решение}
В начале записи указано $\cos^2(3x)=\dfrac35$ (в условии стоит $\dfrac34$). Далее Вы переходите к равенствам
\[
\frac{1-\cos6x}{2}=\frac35,
\qquad 5-5\cos6x=6,
\qquad \cos6x=-\frac15.
\]
В конце написаны две серии корней
\[
x=\frac\pi{18}+\frac{\pi n}{3},\qquad
x=-\frac\pi{18}+\frac{\pi n}{3},\quad n\in\mathbb Z.
\]
\heading{Ваш ответ}
$\displaystyle x=\pm\frac\pi{18}+\frac{\pi n}{3}$, $n\in\mathbb Z$.
\heading{Ошибка}
\err{} В исходном равенстве записана другая правая часть; затем использована формула для квадрата синуса вместо квадрата косинуса. Заключительная запись корней совпадает с правильной, но из предыдущих вычислений она не следует.
\heading{Где именно возникает ошибка}
Первое расхождение с условием появляется ещё при его переписывании:
\[
\underbrace{\cos^2(3x)=\frac35}_{\text{в Вашей записи}}
\qquad\text{вместо}\qquad
\underbrace{\cos^2(3x)=\frac34}_{\text{в условии}}.
\]
Следующая ошибка состоит в знаке перед косинусом двойного угла:
\[
\frac{1-\cos6x}{2}\quad\text{вместо}\quad
\frac{1+\cos6x}{2}.
\]
Даже если бы промежуточное уравнение $\cos6x=-1/5$ было верным, числа $x=\pm\pi/18$ нельзя было бы из него получить, потому что при $x=\pi/18$ косинус $6x$ равен $1/2$, а не $-1/5$.
\heading{Почему это неверно}
Из формулы двойного угла $\cos2a=2\cos^2a-1$ следует
\[
2\cos^2a=1+\cos2a,
\qquad
\boxed{\cos^2a=\frac{1+\cos2a}{2}}.
\]
Знак «минус» относится к другой формуле:
\[
\sin^2a=\frac{1-\cos2a}{2}.
\]
Замена плюса минусом меняет уравнение. Ошибка в правой части дополнительно меняет численное значение косинуса. То, что итоговые серии случайно совпали с верным ответом, не делает эти преобразования правильными.
\heading{Ключевая идея}
\idea{} Сначала точно переносим данные условия, затем выбираем формулу понижения степени в зависимости от того, квадрат \emph{синуса} или \emph{косинуса} записан в уравнении. После нахождения косинуса двойного угла выписываем все серии решений и проверяем их подстановкой.
\heading{Исправление от последнего верного шага}
Возвращаемся к исходному уравнению с правой частью $3/4$ и применяем правильную формулу:
\[
\frac{1+\cos6x}{2}=\frac34
\quad\Longrightarrow\quad
\cos6x=\frac12.
\]
\heading{Полное правильное решение}
Сначала понижаем степень:
\begin{align*}
\cos^2(3x)&=\frac34,\\
\frac{1+\cos(2\cdot3x)}2&=\frac34,\\
1+\cos6x&=\frac32,\\
\cos6x&=\frac12.
\end{align*}
Поскольку $\cos(\pi/3)=1/2$, все решения последнего уравнения таковы:
\[
6x=\pm\frac\pi3+2\pi n,\qquad n\in\mathbb Z.
\]
Делим \emph{каждое} слагаемое на шесть:
\[
\boxed{x=\pm\frac\pi{18}+\frac{\pi n}{3}},
\qquad n\in\mathbb Z.
\]
Знак «плюс-минус» обозначает две серии, обе из которых требуется сохранить.
\heading{Проверка}
Для $x=\pi/18$ имеем $3x=\pi/6$, поэтому $\cos^2(3x)=3/4$. При добавлении $\pi n/3$ к $x$ угол $3x$ увеличивается на $\pi n$, а квадрат косинуса не меняется. Для второй серии $x=-\pi/18+\pi n/3$ проверка аналогична, поскольку косинус чётен.
\heading{Итог}
\result{} Верный ответ: $\displaystyle x=\pm\frac\pi{18}+\frac{\pi n}{3}$, $n\in\mathbb Z$. Он совпадает с Вашим ответом, но путь к нему нужно исправить.
\heading{Задача на закрепление}
Решите уравнение $\displaystyle\cos^2(4x)=\frac14$, используя понижение степени.
\newpage
\phantomsection\label{task:9}
\section*{Задание № 9. Показательное уравнение с тригонометрической степенью}
\heading{Текст задания}
Решите уравнение
\[
4^{2+\cos^2x}-4^{\cos^2x}
=9^{1+\cos^2x}+9^{\cos^2x}.
\]
\heading{Ваше решение}
Подстановка $t=\cos^2x$ и вынесение общего множителя проведены правильно:
\[
4^t(16-1)=9^t(9+1),\qquad
\left(\frac49\right)^t=\frac{10}{15},
\qquad t=\frac12.
\]
Далее записано
\[
\cos^2x=\frac{1-\cos2x}{2}=\frac12,
\qquad \cos2x=0,
\]
и затем указан ответ $x=\pi/8+\pi n/4$.
\heading{Ваш ответ}
$x=\dfrac\pi8+\dfrac{\pi n}{4}$, $n\in\mathbb Z$.
\heading{Ошибка}
\err{} Неверно указан знак в формуле квадрата косинуса; при решении $\cos2x=0$ дополнительно потерян множитель два.
\heading{Где именно возникает ошибка}
До равенства $t=1/2$ преобразования корректны. Первый неверный шаг:
\[
\cos^2x\ne\frac{1-\cos2x}{2}
\quad\text{(в общем случае)}.
\]
Верная формула содержит знак «плюс». В этой конкретной задаче неверный знак не изменил промежуточное равенство $\cos2x=0$, поскольку справа получилось $1/2$. Но следующий переход уже меняет ответ: из $\cos2x=0$ следует $x=\pi/4+\pi n/2$, а не $x=\pi/8+\pi n/4$.
\heading{Почему это неверно}
Формула $\cos2x=2\cos^2x-1$ равносильна $\cos^2x=(1+\cos2x)/2$. Кроме того, уравнение $\cos y=0$ имеет решения $y=\pi/2+\pi n$. Если $y=2x$, нужно сначала приравнять \emph{весь} аргумент $2x$ к $\pi/2+\pi n$, а потом разделить на два ровно один раз.
\heading{Ключевая идея}
\idea{} Разделяйте работу на два этапа: сначала решайте показательное уравнение относительно $t$, затем независимо и аккуратно возвращайтесь к тригонометрическому равенству $\cos^2x=t$.
\heading{Исправление от последнего верного шага}
После корректно полученного $t=1/2$:
\[
\cos^2x=\frac12
\quad\Longrightarrow\quad
\frac{1+\cos2x}{2}=\frac12
\quad\Longrightarrow\quad
\cos2x=0.
\]
Значит,
\[
2x=\frac\pi2+\pi n,
\quad\boxed{x=\frac\pi4+\frac{\pi n}{2}},\quad\nZ.
\]
\heading{Полное правильное решение}
Пусть $t=\cos^2x$. Тогда обязательно $0\le t\le1$. Исходное уравнение принимает вид
\[
4^{2+t}-4^t=9^{1+t}+9^t.
\]
Раскладываем степени и выносим общие множители:
\begin{align*}
16\cdot4^t-4^t&=9\cdot9^t+9^t,\\
15\cdot4^t&=10\cdot9^t,\\
\left(\frac49\right)^t&=\frac23,\\
\left(\left(\frac23\right)^2\right)^t&=\frac23,\\
\left(\frac23\right)^{2t}&=\left(\frac23\right)^1.
\end{align*}
Основание $2/3$ положительно и не равно единице; следовательно, показатель степени однозначен: $2t=1$, откуда $t=1/2$. Найденное значение принадлежит промежутку $[0,1]$, поэтому оно допускается подстановкой $t=\cos^2x$.

Теперь решаем
\[
\cos^2x=\frac12
\quad\Longleftrightarrow\quad
\cos2x=0
\quad\Longleftrightarrow\quad
2x=\frac\pi2+\pi n.
\]
Получаем $\displaystyle x=\frac\pi4+\frac{\pi n}{2}$, $n\in\mathbb Z$.
\heading{Проверка}
Для $x=\pi/4+\pi n/2$ квадрат косинуса равен $1/2$. Тогда левая часть исходного равенства составляет
\[
4^{5/2}-4^{1/2}=32-2=30,
\]
а правая --- $9^{3/2}+9^{1/2}=27+3=30$. Равенство выполняется.
Ваше число $x=\pi/8$ не является корнем, поскольку $\cos^2(\pi/8)=(2+\sqrt2)/4\ne1/2$.
\heading{Итог}
\result{} $\displaystyle x=\frac\pi4+\frac{\pi n}{2}$, $n\in\mathbb Z$.
\heading{Задача на закрепление}
Решите уравнение
\[
16^{1+\sin^2x}-16^{\sin^2x}=30\cdot4^{\sin^2x}.
\]
\newpage
\phantomsection\label{task:10}
\section*{Задание № 10. Уравнение с синусом и косинусом}
\heading{Текст задания}
Решите уравнение
\[
\sin x+\cos x=1.
\]
Объясните, почему прямое деление обеих частей на $\cos x$ требует отдельной проверки случая $\cos x=0$.
\heading{Ваше решение}
Вы верно выразили тригонометрические функции через половинный угол, привели выражение к нулю и вынесли общий множитель:
\[
2\sin\frac x2\cos\frac x2-2\sin^2\frac x2=0,
\]
\[
2\sin\frac x2\left(\cos\frac x2-\sin\frac x2\right)=0.
\]
Получены два случая: $\sin(x/2)=0$ и $\cos(x/2)-\sin(x/2)=0$. Для второго случая после деления на $\cos(x/2)$ записано
\[
1-\tg\frac x2=0
\quad\Longrightarrow\quad
\tg\frac x2=-1.
\]
Несмотря на этот неверный переход, в конце приведена серия $x=\pi/2+2\pi n$. Первая серия в записи: $x=2\pi n$.
\heading{Ваш ответ}
$\displaystyle x=2\pi n$ или $\displaystyle x=\frac\pi2+2\pi n$, $n\in\mathbb Z$.
\heading{Ошибка}
\err{} Из равенства $1-\tg(x/2)=0$ по ошибке получено $\tg(x/2)=-1$ вместо $+1$. При этом Ваши конечные серии корней \emph{верны}; ошибка содержится в обосновании второго семейства.
\heading{Где именно возникает ошибка}
Последний верный шаг:
\[
1-\tg\frac x2=0.
\]
Первый неверный шаг:
\[
\tg\frac x2=-1.
\]
После переноса отрицательного слагаемого в правую часть знак меняется:
\[
1-\tg\frac x2=0
\quad\Longleftrightarrow\quad
\boxed{\tg\frac x2=1}.
\]
\heading{Почему это неверно}
Если $1-a=0$, то $1=a$, следовательно, $a=1$, а не $-1$. При значении $a=-1$ левая часть равенства $1-a$ составляет $2$, то есть нулю не равна. Такое изменение знака ведёт к другой серии углов. В данном случае Вы всё же записали в конце правильный ответ, поэтому исправить требуется именно алгебраический переход.
\heading{Ключевая идея}
\idea{} После разложения произведения на множители приравниваем каждый множитель к нулю. Любое деление на тригонометрическое выражение допустимо только если отдельно доказано, что делитель не обращается в ноль, либо соответствующие нулевые значения уже проверены.
\heading{Исправление от последнего верного шага}
Сохраняем верное разложение:
\[
2\sin\frac x2\left(\cos\frac x2-\sin\frac x2\right)=0.
\]
Если $\sin(x/2)=0$, то $x/2=\pi n$ и $x=2\pi n$. Если же
\[
\cos\frac x2-\sin\frac x2=0,
\]
то $\cos(x/2)$ здесь не равен нулю: иначе одновременно должен был бы равняться нулю $\sin(x/2)$, что невозможно. Значит, можно разделить на $\cos(x/2)$ и получить
\[
\tg\frac x2=1
\quad\Longrightarrow\quad
\frac x2=\frac\pi4+\pi n
\quad\Longrightarrow\quad
x=\frac\pi2+2\pi n.
\]
\heading{Полное правильное решение}
Преобразуем исходное уравнение, используя тождества
$\sin x=2\sin(x/2)\cos(x/2)$ и
$\cos x=1-2\sin^2(x/2)$:
\begin{align*}
\sin x+\cos x&=1,\\
2\sin\frac x2\cos\frac x2+1-2\sin^2\frac x2&=1,\\
2\sin\frac x2\cos\frac x2-2\sin^2\frac x2&=0,\\
2\sin\frac x2\left(\cos\frac x2-\sin\frac x2\right)&=0.
\end{align*}
Произведение равно нулю тогда и только тогда, когда хотя бы один множитель равен нулю.

\textbf{Первый случай:} $\sin(x/2)=0$; отсюда $x/2=\pi n$ и $x=2\pi n$.

\textbf{Второй случай:} $\cos(x/2)=\sin(x/2)$. После обоснованного деления на ненулевой косинус получаем $\tg(x/2)=1$. Значит,
\[
\frac x2=\frac\pi4+\pi n
\quad\Longrightarrow\quad
x=\frac\pi2+2\pi n,\qquad n\in\mathbb Z.
\]
Таким образом,
\[
\boxed{x=2\pi n\quad\text{или}\quad x=\frac\pi2+2\pi n},
\quad n\in\mathbb Z.
\]

\textbf{Почему нельзя сразу делить на $\cos x$.} При делении равенства на выражение $\cos x$ молчаливо предполагают, что оно не равно нулю. Но корни уравнения могут удовлетворять $\cos x=0$. Обязательно нужно проверить этот случай \emph{до} деления:
\[
\cos x=0\ \Longrightarrow\ x=\frac\pi2+\pi n.
\]
Тогда исходное уравнение превращается в $\sin x=1$, откуда
\[
x=\frac\pi2+2\pi n,\qquad n\in\mathbb Z.
\]
Получена целая серия корней, которая при прямом делении на $\cos x$ потерялась бы. После отдельной проверки этого случая деление в оставшемся случае $\cos x\ne0$ допустимо. Заметим, что использованный выше метод разложения на множители вообще не требует деления на $\cos x$.
\heading{Проверка}
При $x=2\pi n$ левая часть равна $0+1=1$. При $x=\pi/2+2\pi n$ она равна $1+0=1$. Других корней нет, потому что два множителя в полученном разложении перебраны полностью.
\heading{Итог}
\result{} $\displaystyle x=2\pi n$ или $\displaystyle x=\frac\pi2+2\pi n$, $n\in\mathbb Z$. Итоговый ответ в тетради верен, однако неверный знак в промежуточном равенстве следует исправить.
\heading{Задача на закрепление}
Решите уравнение $\sin(2x)+\cos(2x)=1$ и укажите, почему нельзя без отдельной проверки делить на $\cos(2x)$.
\vfill
\begin{center}\small Лёвин Артём Александрович эксклюзивно для Николь Саркисьянц\end{center}
\end{document}